\documentclass[12pt]{article}
\usepackage[T1]{fontenc}
\usepackage[utf8]{inputenc}
\usepackage{lmodern}
\usepackage{textcomp}
\usepackage{enumitem}
\usepackage{geometry}
\geometry{margin=1in}
\begin{document}
\begin{center}
Answer Key - Science - K-12 Level Assessment 17 \
\small (Answers are ordered exactly as in the question paper.)
\end{center}

\section*{Multiple Choice}
\begin{enumerate}[label=\arabic*.]
\item \textbf{Answer:} C
\\ \textit{Options:} A) deer; B) raccoons; C) alligators; D) coyotes
\\ \textit{Note:} Alligators are aquatic reptiles that thrive in wetland environments, making them well-adapted to benefit from the increased water levels and habitat expansion that floods can provide in coastal areas.
\item \textbf{Answer:} D
\\ \textit{Options:} A) Top soil loses all of its nutrients.; B) Plants begin to grow deeper roots.; C) Animals have more food sources.; D) Trees become less healthy over time.
\\ \textit{Note:} Acid rain lowers the pH of soil and water, leading to nutrient depletion and harm to trees, which ultimately results in their declining health over time.
\item \textbf{Answer:} B
\\ \textit{Options:} A) They change chemical energy into kinetic energy.; B) They change solar energy into chemical energy.; C) They change wind energy into heat energy.; D) They change mechanical energy into solar energy.
\\ \textit{Note:} Giant redwood trees perform photosynthesis, a process that converts solar energy from sunlight into chemical energy stored in glucose molecules, making the correct answer "They change solar energy into chemical energy."
\item \textbf{Answer:} A
\\ \textit{Options:} A) changing environments.; B) volcanic activity.; C) changing rates of erosion.; D) evaporation of sea water.
\\ \textit{Note:} The lack of fossils in one layer of sedimentary rock is most likely due to changing environments, as different conditions can affect the preservation of organisms and their ability to leave behind fossil evidence.
\item \textbf{Answer:} D
\\ \textit{Options:} A) their preference for eating fish; B) their ability to fly; C) their method of reproduction; D) their method of catching food
\\ \textit{Note:} The correct answer is based on the distinct hunting techniques of eagles, which typically use their keen eyesight and powerful talons to catch prey in mid-air, while pelicans primarily rely on their large beaks to scoop fish from the water, illustrating their different adaptations for obtaining food.
\end{enumerate}

\section*{True / False}
\begin{enumerate}[label=\arabic*.]
\setcounter{enumi}{5}
\item \textbf{Answer:} False
\\ \textit{Options:} A) True; B) False
\\ \textit{Note:} The statement is false because mass is a measure of the amount of matter in an object, which depends on its density and volume; thus, even if the feathers and iron nails occupy the same volume, their different densities mean they have different masses.
\item \textbf{Answer:} False
\\ \textit{Options:} A) True; B) False
\\ \textit{Note:} Grass primarily obtains its energy through the process of photosynthesis by converting sunlight into chemical energy, rather than relying mainly on nutrients from the soil.
\item \textbf{Answer:} False
\\ \textit{Options:} A) True; B) False
\\ \textit{Note:} Moons are defined as natural satellites that orbit planets, and currently, no moons have been confirmed to exist outside our solar system, making the statement false.
\item \textbf{Answer:} True
\\ \textit{Options:} A) True; B) False
\\ \textit{Note:} A star is considered a natural source of light because it generates energy through nuclear fusion in its core, where hydrogen atoms combine to form helium, releasing vast amounts of light and heat in the process.
\item \textbf{Answer:} True
\\ \textit{Options:} A) True; B) False
\\ \textit{Note:} The statement is true because the digestive system is composed of layers of smooth muscle that contract in a coordinated manner to facilitate the movement of food, nutrients, and waste through the gastrointestinal tract.
\end{enumerate}

\section*{Long Form Response}
\begin{enumerate}[label=\arabic*.]
\setcounter{enumi}{10}
\item \textbf{Answer:} The decomposition rate of various materials in nature is influenced by factors such as temperature, moisture, oxygen availability, and the type of material itself. Organic materials like food scraps and plant matter decompose the fastest due to their high nutrient content and the presence of microorganisms that thrive in warm, moist environments. In contrast, materials like plastic and metal take much longer to break down because they are not easily degraded by natural processes. For example, a banana peel can decompose in just a few weeks, while a plastic bottle may take hundreds of years. Overall, the combination of biological activity and environmental conditions plays a critical role in how quickly different materials decompose.
\\ \textit{Note:} The answer correctly identifies that temperature, moisture, oxygen availability, and material type are key factors influencing decomposition rates, and it highlights that organic materials decompose faster due to their nutrient content and favorable conditions for microorganisms, while synthetic materials like plastic and metal decompose much more slowly.
\item \textbf{Answer:} The transformation of sedimentary rock into metamorphic rock occurs through a process called metamorphism, which involves heat, pressure, and mineral changes. When sedimentary rocks, such as sandstone or limestone, are subjected to high temperatures and intense pressure, typically due to tectonic movements or burial deep within the Earth, their minerals can reorganize and alter their structure. This process requires specific conditions, including elevated temperatures of at least 150 to 200 degrees Celsius and significant pressure, often found in subduction zones or mountain-building regions. As a result, the original sedimentary rock can change into metamorphic rock, such as quartzite or marble, which exhibit new physical properties and mineral compositions.
\\ \textit{Note:} This answer is correct because it accurately describes the metamorphic process, emphasizing the roles of heat, pressure, and mineral changes in transforming sedimentary rocks into metamorphic rocks under specific geological conditions.
\end{enumerate}

\vfill
\noindent\textit{End of Answer Key}
\end{document}